\subsection{共享内存中的输入复用}
\label{l11:subsec:shared_memory}

相邻输出的输入窗口大量重叠。对于一个完整的 $T\times T$ 输出分块，固定一个输入通道，如果每个线程都从全局内存独立读取自己的 $K\times K$ 窗口，共发出 $T^2K^2$ 次输入读取请求。其中，同一个输入元素可能被块内多个线程重复请求。

求这些窗口的联合覆盖范围，就能确定需要协作加载的数据。设输出块左上角为 $(h_0,w_0)$，一个线程需要的输入行是 $h_0+t_y+p$，其中 $0\le t_y<T$、$0\le p<K$。它的最小值为 $h_0$，最大值为 $h_0+T+K-2$，因此联合输入区域在高度方向包含 $T+K-1$ 个位置。宽度方向相同，所以输入分块大小为 $(T+K-1)\times(T+K-1)$。

输出分块之外、为完成边缘输出而额外加载的输入区域称为晕区（halo）。将整个联合输入区域加载到共享内存，每个输入元素在这个块内只从全局内存取一次，后面的所有局部乘加都可读取共享缓冲区。对于单个输入通道，完整分块的读取请求减少倍数为
\begin{equation}
  R_{\mathrm{load}}
  =\frac{\text{独立读取各窗口的次数}}{\text{协作加载联合输入区域的次数}}
  =\frac{T^2K^2}{(T+K-1)^2}.
  \label{l11:eq:tile_reuse}
\end{equation}
例如，$T=16,K=7$ 时，独立读取需要 $256\times49=12\,544$ 次输入请求；联合输入块只有 $22^2=484$ 个元素，于是 $R_{\mathrm{load}}=12\,544/484\approx25.92$。若逐个处理 $C$ 个输入通道，分子和分母都乘 $C$，该比值保持不变。

这个比值刻画指定条件下请求的全局输入读取数。共享内存加载、同步、乘加和输出写回仍然需要时间，缓存也可能已服务了基础实现的一部分重复读取，因此不能把 $25.92$ 直接当作运行时间的加速倍数。相邻线程块的晕区仍可能重叠，同一输入区域也可能由不同输出通道对应的线程块重复读取；式\eqref{l11:eq:tile_reuse}只计算一个块内部的协作收益。

\subsection{协作加载、同步与权重复用}
\label{l11:subsec:cooperation}

输入分块有 $(T+K-1)^2$ 个元素，往往多于块内的 $T^2$ 个线程，因此每个线程只装入一个元素并不足够。令 $Q=T+K-1$，块内线性线程号为 $u=t_yT+t_x$。让线程 $u$ 负责所有满足 $a=u+jT^2<Q^2$ 的输入线性位置，其中 $j=0,1,2,\ldots$。每个合法 $a$ 恰好具有一个模 $T^2$ 的余数，因此所有线程的加载范围不重叠，并且覆盖整个输入分块。由 $a$ 恢复共享缓冲区坐标为 $r=\left\lfloor a/Q\right\rfloor$、$s=a\bmod Q$，对应的原输入位置是 $(h_0+r,w_0+s)$。

使用同一共享缓冲区依次处理不同输入通道时，每个通道需要遵守三个阶段的时序：

\begin{enumerate}
  \item 所有线程协作把该通道所需的 $Q\times Q$ 输入区域放入共享内存。边缘块中，超出原输入的加载位置填入 $0$；合法输出所需要的输入仍全部来自有效窗口。
  \item 全块执行同步。随后，负责合法输出的线程读取共享缓冲区中的 $K\times K$ 窗口，把本通道贡献累加到各自的寄存器中。
  \item 在下一通道覆盖这块共享缓冲区之前再次全块同步，确保当前通道的数据已经被所有需要它的线程读完。
\end{enumerate}

第一道同步防止读取尚未加载完的数据，第二道同步防止较快线程过早覆盖较慢线程仍在使用的数据。输出范围外的线程虽然不产生结果，仍可能承担晕区加载并参与同步。因此，代码\ref{l11:lst:conv_gpu}中的越界提前返回只适用于当前没有块内同步的基础核；引入共享缓冲区以后，应保留全块加载与同步，只对累加和最终输出写回设置有效性条件。

按通道复用缓冲区时，共享内存大小约为 $4Q^2$ 字节，而不是 $4CQ^2$ 字节；累加变量跨通道保存在每个线程的寄存器中。一次装入所有通道也能实现相同计算，但共享内存占用随 $C$ 增长。如何选择取决于可用容量、同步成本以及一个计算单元上能够同时驻留多少个线程块。

权重还有另一种复用方式。在固定输出通道的基础映射中，同一线程束（warp）内的线程处理不同空间位置，却在相同的 $c,p,q$ 循环迭代读取相同的 $F_{m,c,p,q}$。当核组能够放入常量内存（constant memory）且访问保持这种同址特征时，可以利用广播让多个线程共享一次权重取值。若另一种线程映射让同一线程束同时处理不同输出通道，权重地址就可能不同，需要重新分析访问行为。代码\ref{l11:lst:conv_gpu}使用普通权重指针，常量内存属于可在该基础上采用的存储优化。

在图像输入方面，固定的 $c,p,q$ 下，相邻横向线程读取相邻列，行主序布局因而有利于连续访问。不同输出通道的线程块还会访问同一幅图像中的相同区域；这些块有可能利用缓存中的输入，但实际复用取决于调度和缓存容量。若希望在一次协作计算中显式共享跨输出通道的输入，可以把多个输出通道一起纳入分块，或者利用下面的矩阵乘法表示。
